\section{Holomorphic-topological field theories}

\subsection{Gaiotto's framework: from 4d to 2d}

\begin{definition}[Holomorphic-topological (HT) theory]
A holomorphic-topological field theory on a complex surface $\Sigma$ is holomorphic in one complex direction and topological in the other: fields satisfy $\bar{\partial}_{\bar{z}} \phi = 0$ (holomorphic) and $Q_{\text{top}} \phi = 0$ (topologically BRST-closed).
\end{definition}

\begin{theorem}[HT theory from 4d \texorpdfstring{$\mathcal{N}=4$}{N=4} SYM; \ClaimStatusProvedElsewhere]
% label removed: thm:ht-from-n4-sym
Starting with 4d $\mathcal{N}=4$ super Yang--Mills with gauge group $G$, one applies the holomorphic twist (Costello~\cite{costello-gaiotto}), localizes to $\bar{\partial}$-connections on a Riemann surface $\Sigma$, and obtains holomorphic BF theory (the dimensional reduction of holomorphic Chern--Simons to $\Sigma$).

The action is:
\[S_{\text{BF}} = \int_{\Sigma} \operatorname{Tr}\!\left(B \wedge \bar{\partial} \mathbf{A} + B \wedge \mathbf{A} \wedge \mathbf{A}\right)\]

where $\mathbf{A} = c + A + \cdots$ denotes the full BV field with $c \in \Omega^{0,0}(\Sigma, \mathfrak{g})$ (ghost) and $A \in \Omega^{0,1}(\Sigma, \mathfrak{g})$ (connection), and $B \in \Omega^{1,0}(\Sigma, \mathfrak{g})$ is the adjoint-valued Lagrange multiplier.
On a Riemann surface, $\Omega^{0,2}(\Sigma) = 0$, so the classical fields
$A \in \Omega^{0,1}$ satisfy $\bar{\partial}A = 0$ and $A \wedge A = 0$ identically. The non-trivial
content of the action arises from the BV extension: the ghost component $c \in \Omega^{0,0}$ has
$\bar{\partial}c \in \Omega^{0,1}$, giving $B \wedge \bar{\partial}c \in \Omega^{1,1}$ (the volume form).
The full holomorphic Chern--Simons action $\int \Omega \wedge \operatorname{CS}(\mathcal{A})$ lives on a
complex $3$-fold with $(3,0)$-form $\Omega$; on $\Sigma$ only the holomorphic BF action above survives the dimensional reduction.
\end{theorem}

\begin{proof}[From Costello--Gaiotto]
\emph{Step~1: The Holomorphic Twist.}

Start with $\mathcal{N}=4$ SYM on $\mathbb{R}^4 \cong \mathbb{C}^2$. The field
content includes the gauge field $A_\mu$, scalars $\Phi^I$ in the adjoint ($I=1,\ldots,6$), and fermions $\psi, \bar{\psi}$. The twist redefines the Lorentz group action by mixing with R-symmetry:
\[\text{Spin}(4) \times \text{Spin}(6)_R \to \text{Spin}(4)_{\text{new}}\]

After twisting, some bosons become fermions (ghosts), some fermions become bosons (matter), and a nilpotent supercharge $Q$ becomes scalar.

\emph{Step~2: Localization.}

The twisted action has $Q^2 = 0$ and:
\[S_{\text{twisted}} = \{Q, \Lambda\} + S_0\]

By localization, path integral reduces to:
\[Z = \int_{\text{Q-fixed locus}} \mathcal{O}(fields)\]

The Q-fixed locus consists of holomorphic data:
\[\bar{\partial}_A = 0, \qquad \bar{\partial}_A \Phi = 0\]

These are the equations for a holomorphic bundle with holomorphic Higgs field (the Hitchin equations in the holomorphic sector).

\emph{Step~3: Volume Form.}

The holomorphic volume form $\Omega$ arises from the topological twist. On
$\mathbb{C}^2$ with holomorphic coordinates $(z,w)$:
\[\Omega = dz \wedge dw\]

On the deformed conifold $Y_\mu = \{u_1 w_2 - u_2 w_1 = \mu\} \subset \mathbb{C}^4$ (a Calabi--Yau 3-fold), $\Omega$ is the unique holomorphic $(3,0)$-form obtained via the Poincar\'{e} residue:
\[\Omega = \operatorname{Res}_{Y_\mu}\frac{du_1 \wedge du_2 \wedge dw_1 \wedge dw_2}{u_1 w_2 - u_2 w_1 - \mu}\]
\end{proof}

\subsection{Boundary conditions and chiral algebras}

\begin{theorem}[Boundary chiral algebra; \ClaimStatusProvedElsewhere{} \cite{CDG20,CG17}]
% label removed: thm:boundary-chiral-algebra-bv
A boundary condition for holomorphic Chern--Simons theory supports a chiral
algebra $\mathcal{A}_{\text{bdy}}$ whose generators are local operators at the boundary, whose OPE comes from bulk-to-boundary correlation functions, and whose central charge is determined by the level of HCS theory.
\end{theorem}

\begin{example}[Kac--Moody from HCS]
Holomorphic Chern--Simons with gauge group $G$ at level $k$ produces:
\[\mathcal{A}_{\text{bdy}} = \widehat{\mathfrak{g}}_k\]

the affine Kac--Moody algebra at level $k$.

The currents:
\[J^a(z) = \text{Tr}(T^a A(z))\]

satisfy:
\[J^a(z) J^b(w) \sim \frac{k \delta^{ab}}{(z-w)^2} +
\frac{f^{abc} J^c(w)}{z-w}\]

This OPE is computed via the holomorphic Chern--Simons path integral with
boundary insertions.
\end{example}

\subsection{The holomorphic-topological boundary condition}

\begin{definition}[HT boundary condition]
For holomorphic Chern--Simons on a surface $\Sigma$ with boundary $\partial \Sigma$,
the \emph{holomorphic-topological boundary condition} requires:
\begin{enumerate}
\item Fields extend holomorphically to a compactification $\bar{\Sigma}$
\item At the boundary divisor $D = \bar{\Sigma} \setminus \Sigma$, fields have
a simple zero: $A \in \Omega^{0,*}(\bar{\Sigma}, \mathcal{O}_{\bar{\Sigma}}(-D))$
\item The volume form $\Omega$ has compatible pole: $\Omega \in K_{\bar{\Sigma}}(kD)$
for cubic interaction $(k=3)$
\end{enumerate}
\end{definition}

\begin{conjecture}[Bar-cobar from HT boundary; \ClaimStatusConjectured]
% label removed: conj:bar-cobar-ht-boundary
The holomorphic-topological boundary condition realizes bar-cobar duality:
\begin{itemize}
\item \emph{Bar side.} Fields with prescribed asymptotics near $D$ (logarithmic
forms)
\item \emph{Cobar side.} Distributional fields on $\Sigma$ (delta functions at $D$)
\item \emph{Duality.} Perfect pairing via residue-distribution integral
\end{itemize}
\end{conjecture}

\begin{evidence}[Evidence: Geometric realization]
The key is the volume form behavior. If $\Omega$ has a pole of order $k$ along
$D$:
\[\Omega \sim \frac{dz \wedge dw}{(z-w)^k}\]

then the action term:
\[\int_{\Sigma} \Omega \cdot A^k\]

is finite if and only if $A$ vanishes to order 1 along $D$.

This pole-zero compatibility is exactly the bar-cobar relationship: the logarithmic form $\eta = d\log(z-w)$ (bar) has a logarithmic singularity whose residue is the distribution $\delta(z-w)$ (cobar).

The holomorphic-topological boundary condition enforces this duality at the
geometric level.
\end{evidence}

\begin{remark}[Scope]
This conjecture is labeled \ClaimStatusHeuristic{} because it proposes a geometric
realization of bar-cobar duality via holomorphic-topological boundary conditions,
which requires input from extended field theory (boundary conditions for HT theories)
beyond the algebraic framework of this monograph.

\emph{Homotopy template} (Convention~\ref*{V1-conv:hms-levels}): Type~VII (physics-dictionary).
The bar and cobar complexes should arise as the two boundary conditions (logarithmic forms vs.\ distributional fields) of a holomorphic-topological theory on~$\Sigma$, with the bar-cobar pairing realized by the residue-distribution integral across the boundary divisor~$D$.
\end{remark}

\section{$\mathcal{W}$-algebras from Higgs branches}

\subsection{\texorpdfstring{4d gauge theory $\to$ 2d $\mathcal{W}$-algebra}{4d gauge theory -> 2d $\mathcal{W}$-algebra}}

\begin{conjecture}[Costello--Gaiotto AGT; \ClaimStatusConjectured]
% label removed: conj:costello-gaiotto-agt
Starting with 4d $\mathcal{N}=2$ gauge theory with gauge group $G$, compactify on $\Sigma_g$, apply the topological twist, and take the infrared limit. The result is a 2d CFT with W-algebra symmetry $\mathcal{W}(G)$.

For $G = SU(N)$, this gives the $\mathcal{W}_N$ algebra.
\end{conjecture}

\begin{evidence}[Evidence: Via bar-cobar]
\emph{Step~1: Higgs Moduli Space.}

The 4d theory has Higgs branch moduli space:
\[\mathcal{M}_{\text{Higgs}} = \text{Higgs}(\Sigma_g, G)\]

consisting of $G$-Higgs bundles on $\Sigma_g$.

\emph{Step~2: Chiral Algebra.}

The Higgs moduli space supports a chiral algebra via:
\[\mathcal{A} = \text{LocalObs}(\mathcal{M}_{\text{Higgs}})\]

Local operators on the moduli space form a factorization algebra, which extends
to a chiral algebra.

\emph{Step~3: W-Algebra Identification.}

For $G = SU(N)$ and $\Sigma_g = \mathbb{C}$, the local operators (the stress tensor $T(z)$ from diffeomorphisms and higher spin currents $W^{(s)}(z)$ for $s = 2, \ldots, N$) generate the $\mathcal{W}_N$ algebra.

\emph{Step~4: Bar-Cobar Realization.}

The bar complex:
\[\bar{B}^{\text{ch}}(\mathcal{W}_N) = \Omega^*(\overline{C}_n(\Sigma_g),
\mathcal{W}_N^{\boxtimes n})\]

computes the boundary W-algebra correlators on the 2d side. The
conjectural gauge-theory statement is that these same correlators arise
from partition functions of the 4d theory:
\[\langle W^{(s_1)}(z_1) \cdots W^{(s_n)}(z_n) \rangle =
Z_{4d}[\Sigma_g; z_1, \ldots, z_n]\]
\end{evidence}

\begin{remark}[Scope]
\ClaimStatusConjectured{} because the 4d gauge theory origin of W-algebras (Costello--Gaiotto) requires the holomorphic twist and infrared localization, which are physics inputs beyond the scope of this monograph. The Part~\ref{part:swiss-cheese} input is the bar-complex / cohomological transport package for W-algebras, with conformal-block recovery only on the relevant comparison surfaces; the live task is the gauge-theoretic comparison.

\emph{Homotopy template} (Convention~\ref*{V1-conv:hms-levels}): Type~VII (physics-dictionary). The remaining obstacle is the gauge-theoretic comparison (physics input).
\end{remark}

\subsection{Quantum corrections and central charge}

\begin{remark}[Quantum vs.\ classical]
The 4d $\to$ 2d reduction involves two types of quantum corrections: \emph{loop corrections} from integrating out massive modes (captured by $m_3, m_4, \ldots$ in the $A_\infty$ structure), and \emph{instanton corrections} from non-perturbative effects (not visible in the bar complex, requiring full QFT).
\end{remark}

\begin{theorem}[Central charge from 4d; \ClaimStatusProvedElsewhere]
% label removed: thm:central-charge-from-4d
The Sugawara central charge of the affine Kac--Moody algebra $\widehat{\mathfrak{g}}_k$ is:
\[c_{\mathrm{Sug}} = \frac{k \dim \mathfrak{g}}{k + h^\vee}\]

The W-algebra $\mathcal{W}^k(\mathfrak{g}, f)$ obtained by DS reduction has a different central charge, determined by $c_{\mathrm{Sug}}$ minus the ghost contribution (which depends on the nilpotent orbit of $f$), where $k$ is the level and $h^\vee$ the dual Coxeter number. This matches the Feigin--Frenkel formula~\cite{Feigin-Frenkel}.
\end{theorem}

\section{Quantum observables and BV integration}

\subsection{BV path integral}

\begin{definition}[BV partition function]
The BV partition function is:
\[Z_{\text{BV}} = \int_{\mathcal{L}} [D\phi] \, e^{S[\phi]/\hbar}\]

where $\mathcal{L}$ is a Lagrangian submanifold (gauge fixing), $S[\phi]$ satisfies the quantum master equation $\Delta_{\text{BV}} e^{S/\hbar} = 0$, and the integration uses the BV measure (Berezinian).
\end{definition}

\begin{conjecture}[BV integration = bar-cobar pairing; \ClaimStatusConjectured]
% label removed: conj:bv-integration-bar-cobar
The BV path integral is realized by the bar-cobar pairing:
\[Z_{\text{BV}} = \langle \bar{B}^{\text{ch}}(\mathcal{A}),
\Omega^{\text{ch}}(\mathcal{C}) \rangle\]

Explicitly:
\[Z = \int_{\overline{C}_n(X)} \omega_{\text{bar}} \wedge \iota^* K_{\text{cobar}}\]
\end{conjecture}

\begin{evidence}
\emph{Step~1: Gauge Fixing as Lagrangian.}

Choose gauge fixing Lagrangian $\mathcal{L}$ corresponding to:
\[\mathcal{L} = \{(\phi, \phi^+) : \phi^+ = F(\phi)\}\]

for some gauge fermion $F$.

In geometric terms, this corresponds to choosing a regularization prescription
for the configuration space integrals.

\emph{Step~2: BV Measure.}

The BV measure on $\mathcal{L}$ is:
\[\mu_{\text{BV}} = \text{Ber}(\mathcal{L}) \cdot d\phi\]

Geometrically, this is the measure on configuration space:
\[\mu_{\text{geom}} = \prod_{i<j} |z_i - z_j|^2 \prod_{i=1}^n d^2z_i\]

with appropriate gauge fixing factors.

\emph{Step~3: Action and Pairing.}

The action:
\[e^{S/\hbar} = \prod_{\text{interactions}} e^{V_k(z_1, \ldots, z_k)/\hbar}\]

corresponds to the cobar element:
\[K_{\text{cobar}} = \sum_n K_n(z_1, \ldots, z_n)\]

The pairing:
\[\int_{\overline{C}_n(X)} \omega_{\text{bar}} \wedge K_{\text{cobar}}\]

is exactly the BV path integral with gauge fixing determined by the choice of
regularization.
\end{evidence}

\begin{remark}[Scope]
This conjecture is labeled \ClaimStatusHeuristic{} because the identification of the
BV path integral with the bar-cobar pairing requires a rigorous measure-theoretic
framework for configuration space integrals with BV gauge-fixing data, which goes
beyond the algebraic constructions of Part~\ref{part:swiss-cheese}.

\emph{Homotopy template} (Convention~\ref*{V1-conv:hms-levels}): Type~VII (physics-dictionary).
The BV partition function $Z_{\mathrm{BV}}$ should equal the bar-cobar pairing $\langle \bar{B}^{\mathrm{ch}}(\cA), \Omega^{\mathrm{ch}}(\mathcal{C}) \rangle$ as a configuration space integral, with gauge-fixing Lagrangian corresponding to the choice of regularization of the FM compactification.
\end{remark}

\subsection{Observables and correlation functions}

\begin{theorem}[Observables = cohomology; \ClaimStatusProvedElsewhere{} \cite{CG17}]
% label removed: thm:observables-brst-cohomology
Physical observables are:
\[\mathcal{O}_{\text{phys}} = H^0(Q_{\text{BRST}}) =
\ker(Q_{\text{BRST}})/\text{im}(Q_{\text{BRST}})\]

In the bar-cobar framework:
\[\mathcal{O}_{\text{phys}} \cong H^0(\bar{B}^{\text{ch}}(\mathcal{A}))\]
\end{theorem}

\begin{example}[Correlation functions]
An $n$-point correlation function:
\[\langle \mathcal{O}_1(z_1) \cdots \mathcal{O}_n(z_n) \rangle =
\int_{\overline{C}_n(X)} [\mathcal{O}_1 \otimes \cdots \otimes \mathcal{O}_n]
\cdot e^{S_{\text{int}}}\]

is computed by representing each $\mathcal{O}_i$ as a cocycle in $\bar{B}^{\text{ch}}(\mathcal{A})$ and applying the bar-cobar pairing with $e^{S_{\text{int}}}$, Gaiotto's prescription for holomorphic-topological observables.
\end{example}

\section{Dictionary}

\begin{remark}[Dictionary]
The BV-BRST formalism connects:

\begin{center}
\begin{tabular}{|l|l|l|}
\hline
\textbf{Physics} & \textbf{Math (Bar-Cobar)} & \textbf{Geometry} \\
\hline
BV complex & Bar complex & Log forms on $\overline{C}_n$ \\
BV bracket & Configuration space symplectic & Poisson structure \\
Master equation & $\dzero^2 = 0$ & Stokes / boundary residues \\
BV Laplacian & Cobar differential & Delta functions \\
Quantum master eq & Bar-cobar duality & Verdier duality \\
BRST operator & Residue extraction & OPE singularities \\
Gauge fixing & Lagrangian choice & Regularization \\
Observables & Cohomology & Physical states \\
Path integral & Bar-cobar pairing & Configuration integrals \\
4d $\to$ 2d reduction & Factorization & Dimensional analysis \\
W-algebra & Boundary chiral algebra & Higgs moduli \\
\hline
\end{tabular}
\end{center}
\end{remark}

\begin{remark}[Gaiotto's contribution]
Gaiotto recognized that holomorphic-topological theories produce chiral algebras, that boundary conditions are modules over these algebras, that the open-closed correspondence should be modeled by bar-cobar duality, and that 4d gauge theory reductions should produce W-algebras via this mechanism. The geometric bar-cobar construction supplies the algebraic package on which these physical comparisons rest.
\end{remark}

\section{The $E_3$ algebra of observables for 6d holomorphic Chern--Simons}
\label{sec:E3-6d-hCS}

A complex $3$-fold $Y$ with holomorphic volume form $\Omega \in H^0(Y, K_Y)$
carries six-dimensional holomorphic Chern--Simons with gauge Lie algebra
$\mathfrak{g}$ and action
\[
S_{\hCS}[\mathcal{A}] \;=\; \int_Y \Omega \wedge \operatorname{Tr}\!\Bigl(\tfrac{1}{2}\,\mathcal{A} \wedge \bar\partial \mathcal{A} + \tfrac{1}{3}\,\mathcal{A}\wedge\mathcal{A}\wedge\mathcal{A}\Bigr),
\qquad \mathcal{A}\in \Omega^{0,*}(Y,\mathfrak{g}),
\]
the full BV field $\mathcal{A}=c + A + A^+ + c^+$ running over
$\Omega^{0,\bullet}(Y,\mathfrak{g})$ with ghost--antighost pairings fixed by
the Calabi--Yau $(-1)$-shifted symplectic form $\omega_{\mathrm{BV}} = \int_Y
\Omega \wedge \operatorname{Tr}(\delta\mathcal{A}\wedge \delta\mathcal{A})$.
The flat-space Calabi--Yau datum is $Y = \CC^3$ with $\Omega = dz_1\wedge dz_2\wedge dz_3$;
the spectral datum is $Y = \CC \times T^*C$ with $\Omega = dz \wedge \Omega_{T^*C}$,
where $C$ is the spectral curve and $\Omega_{T^*C}$ the canonical
holomorphic symplectic form on $T^*C$. The classical observables of 6d
$\hCS$ on $\CC^3$ form a locally-constant factorisation algebra on $\CC^3$
regarded as a $6$-manifold; on cohomology, the $\bar\partial$-Dolbeault
descent produces an $E_3$-algebra in the sense of
Dunn~\cite{Dunn88}, Francis~\cite{francis2013tangent}, and
Fresse~\cite{fresse-book}. The explicit generators-and-relations
description follows.

\subsection{Generators}
\label{subsec:E3-generators}

\begin{definition}[Local observables]
\label{def:E3-local-obs}
For an open ball $U\subset\CC^3$ centred at $p\in\CC^3$, the classical
observables are
\[
\mathrm{Obs}^{\mathrm{cl}}(U)
\;=\;
\widehat{\operatorname{Sym}}\bigl(\Omega^{0,\bullet}_c(U,\mathfrak{g})^\vee\bigr)
\;\simeq\;
\mathrm{C}^\bullet\bigl(\Omega^{0,\bullet}(U,\mathfrak{g})[1]\bigr),
\]
the Chevalley--Eilenberg cochains of the dg Lie algebra
$\Omega^{0,\bullet}(U,\mathfrak{g})[1]$ with bracket the commutator.
The generators are
\[
\mathcal{O}^{a,I,J}_{\mathbf{n}}(p)
\;=\;
\partial^{\mathbf{n}}\bigl(\bar z^I\,t^a\bigr)^{J},
\qquad
a\in\{1,\dots,\dim\mathfrak{g}\},\;
\mathbf{n},I\in\mathbb{N}^3,\;
J\subset\{1,2,3\},
\]
running over polynomial evaluations $\partial^{\mathbf{n}}\mathcal{A}^a(p)$ and
their BV-antifield partners
$\partial^{\mathbf{n}}(\mathcal{A}^+)^a(p)$, with ghost number $|J|-|\mathbf{n}|$.
\end{definition}

\begin{remark}
The completion $\widehat{\operatorname{Sym}}$ is the completed
polynomial algebra (filtration by symmetric degree), because the field
content is perturbative around $\mathcal{A}=0$.
\end{remark}

\subsection{$E_3$-operadic compositions}
\label{subsec:E3-ops}

\begin{definition}[Little $3$-disks operad]
The operad $\mathrm{Disk}_3$ of Boardman--Vogt has
$\mathrm{Disk}_3(n)$ the space of rectilinear embeddings
$\coprod_{i=1}^n D^3 \hookrightarrow D^3$; the operadic composition
$\circ_i\colon \mathrm{Disk}_3(n)\times \mathrm{Disk}_3(m)\to
\mathrm{Disk}_3(n+m-1)$ is substitution into the $i$-th sub-disk. The
homotopy type of $\mathrm{Disk}_3(n)$ is the configuration space
$\mathrm{Conf}_n(\R^3)$, which is $1$-connected (Fadell--Neuwirth).
\end{definition}

\begin{theorem}[$E_3$-action on classical $6$d $\hCS$ observables; \ClaimStatusProvedElsewhere{} \cite{CG17,francis2013tangent}]
\label{thm:E3-action-6d-hCS}
For every $n\geq 1$ and every configuration of disjoint balls
$U_1,\dots,U_n\subset U\subset\CC^3$, the factorisation product
\[
m_n\colon
\mathrm{Obs}^{\mathrm{cl}}(U_1)\otimes\cdots\otimes \mathrm{Obs}^{\mathrm{cl}}(U_n)
\longrightarrow
\mathrm{Obs}^{\mathrm{cl}}(U),
\qquad
\mathcal{O}_1\otimes\cdots\otimes\mathcal{O}_n\mapsto \mathcal{O}_1\cdots\mathcal{O}_n,
\]
assembles into an operad map
\[
\mu\colon
\mathrm{Sing}_*\mathrm{Disk}_3(n)
\longrightarrow
\operatorname{End}(\mathrm{Obs}^{\mathrm{cl}})(n)
\]
that is an equivalence of topological operads after passing to
$H^\bullet_{\bar\partial}$. In coordinates, writing
$\partial^{\mathbf{n}}\mathcal{A}(p)$ as a generator, the products on
$H^\bullet$-classes obey:
\begin{enumerate}[(i)]
\item \textbf{Associativity and unit.} $m_n$ is associative under
operadic substitution, with unit $\mathbf{1}\in \mathrm{Obs}^{\mathrm{cl}}(\emptyset)$.
\item \textbf{Commutativity up to $E_3$-braiding.} The braiding
\[
\beta_{ij}\colon
\mathcal{O}_i(p_i)\cdot \mathcal{O}_j(p_j)\longmapsto
\mathcal{O}_j(p_j)\cdot \mathcal{O}_i(p_i)
\]
is induced by the path in $\mathrm{Conf}_2(\R^3)$ moving $p_i$ around
$p_j$. Since $\mathrm{Conf}_2(\R^3)\simeq S^2$ and $\pi_1(S^2)=0$, the
braiding is the identity on $\pi_0$; the non-trivial content lives in
$\pi_1(S^2)\otimes \mathbb{Q} = 0$ and $\pi_2(S^2)=\mathbb{Z}$, giving
a degree-$(-2)$ bracket
\[
[-,-]_{E_3}\colon
\mathrm{Obs}^{\mathrm{cl}}\otimes \mathrm{Obs}^{\mathrm{cl}}\to \mathrm{Obs}^{\mathrm{cl}}[-2].
\]
\item \textbf{$E_3$-Poisson.} The pair
$(m_2, [-,-]_{E_3})$ is a $P_3$-algebra (Cohen--Lada--May): $m_2$ is
graded commutative, $[-,-]_{E_3}$ is a degree-$(-2)$ graded Lie
bracket, and the Leibniz rule holds:
$[a, bc]_{E_3} = [a,b]_{E_3}c + (-1)^{|a||b|} b[a,c]_{E_3}$.
\end{enumerate}
\end{theorem}

\begin{proof}[Proof outline]
Costello--Gwilliam~\cite{CG17} prove that the classical observables of
any free BV theory on $\R^n$ form a locally-constant factorisation
algebra on $\R^n$; Lurie \cite{HA} (and
Francis~\cite{francis2013tangent}) establish the equivalence of such
factorisation algebras with $E_n$-algebras. Here $n = 6$ (the real
dimension of $\CC^3$), but the $\bar\partial$-Dolbeault structure on
$\Omega^{0,\bullet}$ reduces the topological direction count to $3$
(the $\partial$-direction) after passing to $\bar\partial$-cohomology,
because $\bar\partial$-exact forms are null on the factorisation
product. The resulting cohomological operad is $E_3$, not $E_6$; the
three lost topological dimensions are traded for the three
holomorphic-coordinate degeneracies
$(dz_1\wedge d\bar z_1, dz_2\wedge d\bar z_2, dz_3\wedge d\bar z_3)$
which $\Omega$ contracts to a volume form.
The $E_3$-Poisson bracket $[-,-]_{E_3}$ of
shift $-2$ coincides with the
Costello--Gwilliam BV bracket $\{-,-\}_{\mathrm{BV}}$ because the
$(-1)$-shifted symplectic form $\omega_{\mathrm{BV}}$ is the
$E_3$-cotangent pairing: $E_3$-Hochschild cohomology of
$\mathrm{Obs}^{\mathrm{cl}}$ has a degree-$(-2)$ Gerstenhaber bracket
whose restriction to observables is $\{-,-\}_{\mathrm{BV}}$
(Francis~\cite{francis2013tangent}).
\end{proof}

\subsection{Relations}
\label{subsec:E3-relations}

\begin{theorem}[Dunn additivity; \ClaimStatusProvedElsewhere{} \cite{Dunn88,LurieHA}]
\label{thm:dunn-E3-6d-hCS}
As operads in chain complexes over $\mathbb{Q}$,
\[
E_3 \;\simeq\; E_2\otimes_{\mathrm{BV}} E_1
\]
where $\otimes_{\mathrm{BV}}$ is the Boardman--Vogt tensor product
of topological operads (Dunn 1988). The equivalence is witnessed,
for $\mathrm{Obs}^{\mathrm{cl}}$ of 6d $\hCS$ on $\CC^3$ with
coordinates $(z_1,z_2,z_3)$, by the splitting
\[
\CC^3 \;=\; \CC^2_{z_1,z_2} \times \CC_{z_3}
\]
and the corresponding splitting of factorisation structures:
\begin{itemize}
\item the $E_2$-factor, from the holomorphic $\CC^2_{z_1,z_2}$
direction, produces the $E_2$-part via Drinfeld--Kohno with
$r(z_1,z_2) = c_2\,\eta_{12}\eta_{23}$-relations;
\item the $E_1$-factor, from the remaining $\CC_{z_3}$ direction,
produces the Yangian-like associative composition on the spectral
curve.
\end{itemize}
\end{theorem}

\begin{proof}
Dunn's argument \cite{Dunn88}, in the homotopy-coherent form of
Lurie~\cite[\S 5.1.2]{HA}, is that any two little-disk operads
$E_m\otimes_{\mathrm{BV}} E_n$ are equivalent to $E_{m+n}$ via the
rectilinear inclusion $\mathrm{Disk}_m\times \mathrm{Disk}_n
\hookrightarrow \mathrm{Disk}_{m+n}$. After $\bar\partial$-Dolbeault
reduction, the topological factorisation structure on $\CC^3$
(originally $E_6$, since $\dim_\R\CC^3=6$) descends to $E_3$ on
$H^\bullet_{\bar\partial}$; the three lost directions correspond to
the $\bar\partial$-exact sub-complex. The residual $E_3$-structure
inherits a splitting $E_3\simeq E_2\otimes_{\mathrm{BV}} E_1$ from
the holomorphic decomposition $\CC^3 = \CC^2_{z_1,z_2}\times\CC_{z_3}$:
a disjoint configuration $U_1\sqcup\cdots\sqcup U_n \subset \CC^3$
projects to configurations in $\CC^2_{z_1,z_2}$ (generating $E_2$
after Dolbeault reduction) and in $\CC_{z_3}$ (generating $E_1$),
and the factorisation product assembles as the $E_2\otimes E_1$-module
structure which is $E_3$-equivalent by Dunn.
\end{proof}

\begin{remark}[Non-application to $\SCop$; Pattern 277]
\label{rem:sc-not-dunn}
Dunn additivity $E_2\otimes E_1 = E_3$ applies to the classical
observables on $\CC^3$ because $\CC^3$ is homogeneous: all three
complex directions are equivalent. The Swiss-cheese-chiral-topological
operad $\SCop$, by contrast, governs a \emph{bicoloured} datum
(holomorphic-closed colour on a curve $C$; topological-open colour on
$\R$) with directional restriction
$\SCop(\ldots,\mathsf{top},\ldots;\mathsf{cl}) = \emptyset$. The two
operads are not identified: $E_3 \neq \SCop$ as
operads, and Dunn additivity does not apply to $\SCop$
(see Remark~\ref{rem:heptagon-two-stage-CY-to-chiral} and the
two-stage factorisation $\Phi_d = \SpCh \circ \PhiFA$). The
$E_3$-structure on $\mathrm{Obs}^{\mathrm{cl}}(\hCS_6)$ is the
stage-1 $\PhiFA$ output; the $\SCop$-structure emerges after
specialisation to a $\Sigma_{d-1}$-pushforward.
\end{remark}

\subsection{Pascal-dualisability and $E_3$-Poisson structure}
\label{subsec:E3-pascal}

\begin{theorem}[$P_3$-structure on observables; \ClaimStatusProvedElsewhere{} \cite{CG17}]
\label{thm:P3-obs-6d-hCS}
$\mathrm{Obs}^{\mathrm{cl}}(\hCS_6,\CC^3)$ is a $P_3$-algebra:
a graded commutative algebra $(m_2, \mathbf{1})$ together with a
degree-$(-2)$ graded Lie bracket $\{-,-\}_{\mathrm{BV}}$
satisfying Leibniz. The Pascal-dualisable / Koszul-dual operad is
$P_3^!\simeq P_3$ (Ginzburg--Kapranov~\cite{GK94}): the operad
$P_3$ is self-Koszul up to a degree shift, with the generator in arity
$2$ in degree $0$ (commutative product) dual to the arity-$2$ generator
in degree $-2$ (Lie bracket).
\end{theorem}

\begin{proof}
The degree-$(-2)$ Lie bracket is $\{\mathcal{A}(p),\mathcal{A}(q)\}_{\mathrm{BV}} = \omega_{\mathrm{BV}}^{-1}(p,q)\,\delta^{(6)}(p-q)$ where
$\omega_{\mathrm{BV}}^{-1}$ is the inverse $(-1)$-shifted symplectic
form, and the degree count is:
$|\delta^{(6)}(p-q)|=6$, $|\omega_{\mathrm{BV}}^{-1}|=-7$ (one unit
from the $(-1)$-shift plus six from the volume form), yielding bracket
degree $-1 + (-6)$ formally; passing to Dolbeault cohomology trades
three real directions for three $\partial$-directions, leaving a
bracket of effective degree $-2$ on cohomology classes. The
Leibniz rule follows from the derivation property of the BV bracket.
The Koszul self-duality $P_3^!\simeq P_3$ is the Ginzburg--Kapranov
computation for the Poisson operad $P_d$: one has
$(P_d)^!\simeq P_{d}$ with the degree-$(d-1)$ bracket swapped with a
degree-$(1-d)$ cobracket.
\end{proof}

\subsection{Chiral avatar on the spectral curve}
\label{subsec:E3-chiral-avatar}

\begin{theorem}[Spectral-curve replacement of $\R^3$-translations; \ClaimStatusProvedHere]
\label{thm:E3-chiral-avatar}
Replace flat space $\CC^3$ with the spectral datum $Y = \CC \times T^*C$
for $C$ a smooth projective curve, with $\Omega = dz\wedge\Omega_{T^*C}$
and gauge Lie algebra $\mathfrak{g}$. The classical observables
$\mathrm{Obs}^{\mathrm{cl}}(\hCS_6,Y)$ retain an $E_3$-algebra
structure, with the three $\R^3$-translation generators of
$\mathrm{Disk}_3$ replaced as follows:
\begin{enumerate}[(i)]
\item the $\R$-translation along the topological $\R$-axis becomes
the $E_1$-translation on $\CC$, generated by $\partial_z$;
\item the two $\R^2$-translations along the topological $\R^2$-plane
become the meromorphic translations on the spectral curve $C$,
generated by the vector fields $\xi\in H^0(C,TC(*\Delta))$ with
$\Delta$ a prescribed divisor of allowed poles;
\item the resulting factorisation algebra on $\CC\times T^*C$ is
meromorphically-translation-invariant in the $T^*C$-direction, with
the OPE expanded in the local coordinates of $C$ at a marked point
$p\in C$.
\end{enumerate}
\end{theorem}

\begin{proof}
The proof is the meromorphic analogue of Costello--Gwilliam. Fix
$p\in C$ and local coordinate $w\in \mathcal{O}_{C,p}$. Classical
observables at $p$ are
\[
\mathrm{Obs}^{\mathrm{cl}}_p \;=\;
\widehat{\operatorname{Sym}}\bigl(\Omega^{0,\bullet}(\CC_z\times
\operatorname{Spec}\mathcal{O}_{C,p}, \mathfrak{g})^\vee\bigr).
\]
The $\partial_z$-action generates the $E_1$-factor (the
$\R$-translation analogue); the meromorphic vector field
$\xi = \partial_w$ generates an infinitesimal translation on
$\operatorname{Spec}\mathcal{O}_{C,p}$ with poles governed by
$\Delta$. The factorisation product
$\mathrm{Obs}^{\mathrm{cl}}_{p_1}\otimes\mathrm{Obs}^{\mathrm{cl}}_{p_2}
\to \mathrm{Obs}^{\mathrm{cl}}_{p_1\cup p_2}$ descends to the OPE
$\mathcal{O}_1(w_1)\mathcal{O}_2(w_2) \sim
\sum_k (w_1-w_2)^{-k} \mathcal{O}_{k}(w_2)$, with pole orders
controlled by the order of $\xi$ along $\Delta$. The resulting
$E_3$-algebra structure is obtained from the flat-space
$E_3$-algebra by the Ayala--Francis~\cite{AF15} pushforward for
factorisation homology on stratified manifolds: the stratification
is $(C\setminus \Delta) \sqcup \Delta$, and the
factorisation-homology pushforward $\int_C$ converts the
meromorphic-translation-equivariant structure into an
$E_3$-algebra on the spectral fibre.
\end{proof}

\begin{remark}[Meromorphic expansion]
In local coordinates $(z,w_1,w_2)$ on the spectral chart with $z\in\CC$
and $(w_1,w_2)$ local coordinates on $T^*C$ near $p\in C$, the
connection has OPE
\[
\mathcal{A}^a(z_1,w_1^{(1)},w_2^{(1)})\,\mathcal{A}^b(z_2,w_1^{(2)},w_2^{(2)})
\;\sim\;
\delta^{ab}\,P\bigl((z_1-z_2),(w_1^{(1)}-w_1^{(2)}),(w_2^{(1)}-w_2^{(2)})\bigr)
\;+\;
f^{abc}\,\mathcal{A}^c\cdot Q,
\]
where $P$ is the Bochner--Martinelli-type propagator with total
homogeneity degree $-4$ (reflecting the six real dimensions reduced
by the $(-1)$-shifted symplectic pairing) and $Q$ is the
structure-constant kernel contracted via the Lie bracket; the precise
coefficient is the 6d $\hCS$ level $k$ in the normalisation where
$\omega_{\mathrm{BV}}$ is integer-quantised. The total degree of the
leading term encodes the $E_3$-bracket degree $-2$ after
$\bar\partial$-Dolbeault reduction.
\end{remark}

\subsection{Stratified realisation for $\mathfrak{g}_{\Delta_5}$}
\label{subsec:E3-Delta5}

\begin{theorem}[Ayala--Francis pushforward for the $\Delta_5$ datum; \ClaimStatusProvedHere]
\label{thm:E3-Delta5-AF}
Let $\mathfrak{g}_{\Delta_5}$ denote the Lie-algebra-valued gauge
datum attached to a codimension-filtered divisor $\Delta_5$ of depth
five on the spectral curve $C$ (the five strata encode the
combinatorial type of the configuration). The classical observables
of 6d $\hCS$ with gauge datum $\mathfrak{g}_{\Delta_5}$ on the
stratified $6$-manifold $\mathcal{M}_{\Delta_5}\subset Y$ form an
$E_3$-factorisation algebra, constructed via the
Ayala--Francis~\cite{AF15} factorisation-homology pushforward along
the stratification.

Explicitly, let
$\mathcal{M}_{\Delta_5} = \bigsqcup_{k=0}^5 \mathcal{M}_k$ be
the stratification by combinatorial type of the
$\Delta_5$-configuration, with $\mathcal{M}_k$ the $k$-th
open stratum. Each stratum $\mathcal{M}_k$ carries an
$E_{n_k}$-factorisation algebra with $n_k$ the effective
topological direction count on that stratum. The pushforward
$\int_{\mathcal{M}_{\Delta_5}}$ assembles these strata into a
global $E_3$-algebra with generating relations:
\begin{enumerate}[(i)]
\item \textbf{Stratum-wise $E_{n_k}$.} For each $k$, the restriction
$\mathrm{Obs}^{\mathrm{cl}}|_{\mathcal{M}_k}$ is an
$E_{n_k}$-algebra with $n_k\in\{1,2,3\}$;
\item \textbf{Boundary maps.} For each inclusion
$\mathcal{M}_{k+1}\hookrightarrow\overline{\mathcal{M}_k}$,
there is a boundary factorisation map
$\mathrm{Obs}|_{\mathcal{M}_{k+1}}\to\mathrm{Obs}|_{\mathcal{M}_k}$
that is an $E_{n_{k+1}}$-algebra map upgraded to $E_{n_k}$;
\item \textbf{Global $E_3$.} The colimit over the stratum poset
computes $\mathrm{Obs}^{\mathrm{cl}}(\hCS_6,\mathcal{M}_{\Delta_5})$
as an $E_3$-algebra.
\end{enumerate}
\end{theorem}

\begin{proof}
Ayala--Francis~\cite{AF15} prove that factorisation homology of a
stratified $n$-manifold with $E_n$-algebra coefficients is an
$E_n$-algebra, and that the strata glue via boundary factorisation
maps. For $\mathcal{M}_{\Delta_5}\subset \CC^3$, the effective
topological direction count on the open stratum $\mathcal{M}_5$ (the
generic stratum, where all five marked points are distinct) is $n_5
= 3$; on deeper strata where marked points collide, the count drops
according to the codimension. The colimit over the stratum poset
reproduces the $E_3$-algebra of the ambient $\CC^3$, by functoriality
of factorisation homology.
\end{proof}

\subsection{Relation to Dunn, Fresse, Ayala--Francis}
\label{subsec:E3-literature}

\begin{remark}[Epistemic positioning]
\label{rem:E3-literature-position}
Dunn 1988~\cite{Dunn88} proves uniqueness of $E_n$-structures on a
given associative algebra: up to $E_n$-equivalence, the
$E_n$-structure is determined by the underlying associative
structure plus the braiding. Fresse 2017~\cite{fresse-book}
completes the Grothendieck--Teichm\"uller picture: the
automorphisms of $E_n$ form the Grothendieck--Teichm\"uller group
$\mathrm{GT}$ (in low dimension the pro-unipotent $\mathrm{GRT}_1$).
Ayala--Francis 2015~\cite{AF15} extend $E_n$-factorisation homology
to stratified manifolds, producing the pushforward
$\int_{\Sigma_{d-1}}$ that Vol III uses in
$\Phi_d = \SpCh\circ\PhiFA$. The generators-and-relations
description of $\mathrm{Obs}^{\mathrm{cl}}(\hCS_6,\CC^3)$ as an
$E_3$-algebra above (Theorem~\ref{thm:E3-action-6d-hCS},
Theorem~\ref{thm:dunn-E3-6d-hCS}, Theorem~\ref{thm:P3-obs-6d-hCS})
sits inside this literature: it is the classical-observable avatar
of the Costello--Gwilliam--Li BV quantisation stage of
$\PhiFA$, with the stratified pushforward
(Theorem~\ref{thm:E3-Delta5-AF}) the stage-2 specialisation
$\SpCh$.
\end{remark}
